\section{Related Work}
\label{sec:related}

\subsection{Graph partitioning algorithms}

The METIS multilevel scheme~\citep{karypis1998} introduced coarsening,
initial partitioning, and uncoarsening with Fiduccia--Mattheyses (FM)
refinement.
Both recursive bisection and direct $k$-way use a single multilevel
coarsening-uncoarsening pass: the graph is coarsened to approximately
$k \cdot c$ nodes (with $c \approx 20$), partitioned, and then
uncoarsened with FM refinement in a single pass.
The critical difference is that recursive bisection applies this
single-pass multilevel scheme once per bisection node---$O(\log k)$ times
in total across the bisection tree---while direct $k$-way applies it once
for the full $k$-way problem.
The direct $k$-way variant therefore avoids the overhead of re-coarsening
at each bisection level, and typically runs faster for large
$k$~\citep{karypis1999multilevel}.

SCOTCH~\citep{pellegrini2008scotch} and KaHyPar~\citep{kahypar2017} also
implement both strategies; the relative quality rankings are broadly
consistent with METIS results, with RB generally producing lower edge cuts
on planar graphs due to better exploitation of separator structure.

Sanders and Schulz~\citep{sanders2013think} showed that local search after
both RB and k-way partitioning can close the quality gap, but that RB
provides a better starting point for planar graphs.

\subsection{Redistricting as graph partitioning}

The graph-partitioning formulation of redistricting originates with
Garfinkel and Nemhauser~\citep{garfinkel1970optimal}, who posed it as an
integer program.  Mehrotra et al.~\citep{mehrotra1998optimal} developed an
optimisation heuristic; Bozkaya et al.~\citep{bozkaya2003zone} applied tabu
search; Bação et al.~\citep{bacao2005applying} used genetic algorithms.
None of these works evaluated the RB vs.\ n-way tradeoff explicitly.

Markov chain Monte Carlo approaches~\citep{deford2019recombination,fifield2020automated}
sidestep the RB/n-way dichotomy by using recombination moves, but scale
poorly beyond $k \sim 20$.  Altman and McDonald~\citep{altman2005contiguity}
showed that contiguity constraints (enforced naturally by both RB and k-way
on connected graphs) have measurable effects on electoral competition.

\subsection{Compactness measures}

Polsby and Popper~\citep{polsby1991third} introduced the isoperimetric
compactness ratio $\mathrm{PP} = 4\pi A / P^2$, where $A$ is area and
$P$ is perimeter.  Reock~\citep{reock1961measuring} proposed a minimum
bounding circle ratio.  We use PP as the primary metric because it directly
corresponds to the edge-cut objective on boundary-length-weighted graphs
(B.2): minimising total boundary perimeter is equivalent to minimising
total edge weight for shared-boundary edge weights.

\subsection{Prior work in this portfolio}

B.3 established that for single-objective congressional maps (unweighted),
RB and n-way are statistically equivalent.  B.4 extended this to
edge-weighted formulations and found a small but consistent compactness
advantage for RB.  The present paper generalises both results to arbitrary
chamber types, provides a runtime characterisation across the full $k$
range, and identifies the prime-$k$ edge case.
